% Archivo GENERADO por 04_analisis/scripts/02_caracterizacion.py — no editar a mano
\begin{tabular}{@{}>{\raggedright\arraybackslash}p{5.2cm}ccccc@{}}
  \toprule
  Parámetro & EST & AS1 & AS2 & AGR1 & AGR2 \\
  \midrule
  pH & \num{7.47} & \num{6.35} & \num{8.07} & \num{6.80} & \num{7.61} \\
  Temperatura (\unit{\degreeCelsius}) & \num{17.8} & \num{16.8} & \num{16.9} & \num{18.1} & \num{18.4} \\
  Conductividad eléctrica (\unit{\micro\siemens\per\centi\meter}) & \num{156} & \num{272} & \num{585} & \num{312} & \num{608} \\
  Potencial de óxido-reducción (\unit{\milli\volt}) & \num{214} & \num{-27.2} & \num{88.6} & \num{34.3} & \num{90.8} \\
  Oxígeno disuelto (\unit{\milli\gram\per\liter}) & \num{6.32} & \num{2.69} & \num{5.17} & \num{3.22} & \num{5.37} \\
  Sólidos disueltos totales (\unit{\milli\gram\per\liter}) & \num{78.0} & \num{137} & \num{292} & \num{156} & \num{304} \\
  Turbiedad (NTU) & \num{20.6} & \num{99.0} & \num{190} & \num{156} & \num{30.8} \\
  DQO total (\unit{\milli\gram\per\liter}) & \num{6.33} & \num{411} & \num{129} & \num{393} & \num{75.7} \\
  DQO disuelta (\unit{\milli\gram\per\liter}) & \num{3.33} & \num{348} & \num{54.0} & \num{309} & \num{62.3} \\
  DBO$_5$ (\unit{\milli\gram\per\liter}) & \num{2.23} & \num{219} & \num{10.8} & \num{104} & \num{8.60} \\
  SST (\unit{\milli\gram\per\liter}) & \num{31.2} & \num{82.0} & \num{406} & \num{137} & \num{595} \\
  Amonio (NH$_4^+$) (\unit{\milli\gram\per\liter}) & \num{0.15} & \num{1.04} & \num{1.82} & \num{12.8} & \num{2.41} \\
  Nitrato (NO$_3^-$) (\unit{\milli\gram\per\liter}) & \num{4.83} & \num{11.0} & \num{4.00} & \num{4.55} & \num{26.2} \\
  Nitrógeno total (\unit{\milli\gram\per\liter}) & \num{0.90}\textsuperscript{(5)} & \num{2.72} & \num{2.00} & \num{11.6}\textsuperscript{(5)} & \num{4.46}\textsuperscript{(5)} \\
  Fósforo total (\unit{\milli\gram\per\liter}) & \num{0.10} & \num{0.36} & \num{5.30} & \num{6.33} & \num{3.92} \\
  Fosfato (PO$_4^{3-}$) (\unit{\milli\gram\per\liter}) & \num{1.22} & \num{2.86} & \num{2.80} & \num{10.4} & \num{1.45} \\
  Coliformes totales (NMP/\qty{100}{\milli\liter}) & \num{1.5e4} & \num{1.4e8} & \num{8.6e6} & \num{7.9e5} & \num{4.3e4} \\
  \textit{E.~coli} (NMP/\qty{100}{\milli\liter}) & \num{1.1e3} & \num{7.6e5} & \num{1.0e5} & \num{1.4e4} & \num{2.2e3} \\
  \bottomrule
  \multicolumn{6}{@{}p{0.97\textwidth}@{}}{\footnotesize Media aritmética de las campañas; coliformes: media geométrica. Fuente de los datos: informes del Laboratorio del Agua (FCI, UCM). Se excluye la campaña~1 de AS1 y AS2 (muestras mal tomadas): $n = 6$ en EST, AGR1 y AGR2, $n = 5$ en AS1 y $n = 1$ en AS2 (las muestras de AS2 desde la campaña~3 no llegaron al laboratorio); entre paréntesis, el $n$ cuando difiere. Nitrato informado como NO$_3^-$; nitrógeno total no medido en la campaña~1; DBO$_5$ de la campaña~4 medida a los 4,7~días. EST: estero Batuco; AS1: oficinas; AS2: casa; AGR1: bodega de agroquímicos; AGR2: sala de bombas.}
\end{tabular}
